% ID: SEL-D0008-BA-E02
% ESTADO: cotejada. Segunda pasada visual completa y borrador compilado revisados.
% Fuente: Selectividad/matematicas_aplicadas/ORD_SUPLENTE_Modelo-4-Enunciados (Mat-Ap-2021).pdf
% SHA-256: 74b067b81be20aa0262a0cc58795f3b25488aa9c749d825b7b33f656dd1ae089
% Página: 1; bloque A; ejercicio original 2.

Se consideran las siguientes inecuaciones:
\[
5x-3y\geq -9\qquad x+y\leq 11\qquad 6x+y\leq 36\qquad x+2y\geq 6
\]
\begin{enumerate}[label=\alph*)]
  \item \textbf{(1.5 puntos)} Represente la región factible definida por las
  inecuaciones anteriores y determine sus vértices.
  \item \textbf{(0.25 puntos)} \textquestiondown{}Pertenece el punto $(5,7)$ a
  la región factible anterior?
  \item \textbf{(0.75 puntos)} Calcule los valores máximo y mínimo de la
  función $F(x,y)=10x-6y$ en la región anterior y determine los puntos en los
  que se alcanzan.
\end{enumerate}
